\section{The Encoder Track}
\label{sec:encoder}

\paragraph{Model and training.}
We fine-tuned all 435M parameters of DeBERTa-v3-large \citep{he2023debertav3} as a
9-way classifier, with the \texttt{[CLS]} state passed through a dense tanh pooler and a
new linear layer over the nine labels. The baseline input pairs the sub-question (at most
96 tokens) with the answer, cut from the right at 512 tokens. The final input puts
``Sub-question: \ldots{} Full question: \ldots'' first, with up to 256 tokens, and raises
the context to 1,024 tokens. We added the full question because 69\% of training rows
share their answer with another sub-question, and a model that cannot see the other
sub-questions has no way to tell which part of the answer responds to which. Training
used plain cross-entropy, AdamW at $10^{-5}$ with layer-wise decay 0.95, a cosine schedule
and batch size 16 (Appendix~\ref{sec:app}), and the epoch was chosen on the train slice.
A seed fixes the head initialisation and the data order, so seeds 0--9 are paired across
configurations.

\begin{table}[t]
\centering
\small
\setlength{\tabcolsep}{4pt}% local to this table: five columns in one column width
\begin{tabular}{@{}lcccc@{}}
\toprule
 & \multicolumn{2}{c}{Single model} & \multicolumn{2}{c}{System} \\
\cmidrule(lr){2-3}\cmidrule(l){4-5}
Configuration & \Stwo & \Sone & \Stwo & \Sone \\
\midrule
Baseline, 8 ep. & \pmsd{0.315}{0.040} & \pmsd{0.576}{0.030} & 0.428 & 0.601 \\
+ 16 epochs & \pmsd{0.362}{0.026} & \pmsd{0.601}{0.016} & 0.400 & 0.620 \\
+ full question & \pmsd{0.384}{0.030} & \pmsd{0.614}{0.027} & 0.405 & 0.648 \\
\midrule
All of train$^{\dagger}$ & \pmsd{0.416}{0.020} & \pmsd{0.614}{0.016} & 0.489 & -- \\
\bottomrule
\end{tabular}
\caption{DeBERTa-v3-large on dev. Single model is the mean\,$\pm$\,sd over seeds 0--9.
System is the 10-seed mean probability with logit adjustment ($\tau$ by nested CV). Rows 2
and 3 each add one change, and row 3 is the final system. $^{\dagger}$All of train, last
epoch, seeds 0--4 only, so its system is a 5-seed number and not comparable.}
\label{tab:enc-main}
\end{table}

\paragraph{What improved the model.}
Two changes improved single models (Table~\ref{tab:enc-main},
Figure~\ref{fig:enc-seeds}). Training for 16 epochs instead of 8 raised \Stwo{} by
+0.047 ($t=3.4$, 9 of 10 seeds up) and \Sone{} by +0.026 ($t=2.7$, 8 of 10). Adding the
full question gave a further +0.022 on \Stwo{} ($t=2.0$, 7 of 10) and +0.012 on \Sone{}
($t=1.3$, 6 of 10). Together the two changes add +0.069 on \Stwo{} and +0.038 on \Sone{},
each on 9 of 10 seeds. Most of the gain came from longer training, which also helped the
classes the baseline handled worst. General, for example, rose from 0.119 to 0.295 per
model (seeds 0--4).

\FigureEncSeeds

At first, undertraining looked like a negative result. With 8 epochs the full question
scored below the baseline (\pmsd{0.282}{0.043} against \pmsd{0.337}{0.044}, seeds 0--4),
but its final training loss (1.32--1.50, against 0.57--1.09) showed it had not finished
fitting. Its seeds were less confident (mean top probability 0.445 against 0.594) and
leaned on the class prior, and at 16 epochs the same input gave the best single model.

\paragraph{From model to system.}
Better single models did not give a better \Stwo{} system. Logit adjustment lifted the
baseline ensemble by +0.109 (to 0.428) but did nothing for the final ensemble (0.409 to
0.405). The \Stwo{} difference of $-$0.022 between the two systems has a 95\% bootstrap
interval of [$-$0.101, +0.061], while on \Sone{} the final system is ahead (0.648 against
0.601). We think, although we have not tested it, that longer training and logit
adjustment fix the same fault, the pull towards frequent classes. We also learnt that
five-seed systems are unreliable on 308 items. The baseline system scored 0.438 on seeds
0--4 and 0.356 on seeds 5--9, a spread larger than most of the effects we wanted to
measure, so all our system claims use ten seeds. Training on all of train for 16 epochs
and keeping the last epoch gained +0.029 on \Stwo{} per model on seeds 0--4 (4 of 5 up)
and lost 0.017 on \Sone{}, which with five seeds is not enough to rank it.

\paragraph{What did not help.}
Seven further ideas, each tested against a named control, left the system where it was.
A second encoder that re-ranked the baseline's top three labels using their definitions
scored 0.354 against 0.365 for the ensemble it re-ranked, which suggests that learning
what each label means from 3,448 examples is harder than classifying them. Balanced
Softmax with focal loss \citep{ren2020balanced,lin2017focal} over-corrected towards rare
classes (Explicit F1 fell from 0.679 to 0.319) and used up the gain that logit adjustment
would otherwise have given. Three hierarchy designs (routing through the clarity level, a
factorised head, and a Non-Reply gate with branch specialists) did not beat the flat
softmax, which already scores all nine labels jointly. Boundary experts for the three
most-confused pairs changed \Stwo{} by only +0.002, because they rarely overturn a decision
the flat model makes from the same input. A uniform soup of the ten models scored 0.267,
since seeds with different head initialisation and data order do not average well. Nine
per-class decision weights overfit the 308 items, and an ensemble mixing both inputs scored
0.403 against 0.405. Since nothing trained on the same 3,448 examples moved the system,
\S\ref{sec:llm} tests the backbone itself.
